\section{Fréquence d'une valeur}

\begin{definitionbox}
    La \bb{fréquence d'une valeur} est le quotient de l'effectif de cette
    valeur par l'effectif total.
    \begin{center}
        $\text{Fréquence} = \dfrac{\text{Effectif}}{\text{Effectif total}}$
    \end{center}
\end{definitionbox}

\begin{exempleboxsuite}
    Pour \bb{calculer la fréquence} des élèves venant à pied au collège on
    connait :

\begin{itemize}[label=\textbullet]
    \item l'\bb{effectif} des élèves venant à pied : \answer{117}
    \item l' \bb{effectif total} : \answer{585}
\end{itemize}

    On en déduit la fréquence des élèves venant à pied : $f$ = 
    \answer{$\dfrac{117}{585} = \dfrac{1}{5} = 0.2 = 20\%$}
\end{exempleboxsuite}

\begin{remarquebox}
    Une fréquence peut se donner
    \begin{itemize}[label=\textbullet]
        \item sous forme d'une fraction simplifiée : $f$ = \answer{$\dfrac{1}{5}$}
        \item en écriture décimale : $f$ = \answer{$0.2$}
        \item en pourcentage : $f$ = \answer{$20\%$}
    \end{itemize}

    \vspace{20pt}

    On peut également présenter les résultats dans un \bb{tableau des
    fréquences} :

    \begin{center}
    \begin{tabular}{lccccc}
       \bb{Transport} & \bb{Deux roues} & \bb{Bus} & \bb{À pied} & \bb{Voiture parents} & \bb{Total}\\
       \midrule
       \bb{Effectif} & \answer{71} & \answer{305} & \answer{117} & \answer{92} & \answer{585} \\
       \bb{Fréquence} & \answer{$\dfrac{71}{585}$} & \answer{$\dfrac{61}{117}$} & \answer{$\dfrac{1}{5}$} & \answer{$\dfrac{92}{585}$} & \answer{$1$} \\
       \bb{Fréquences en \%} & \answer{$12.14\%$} & \answer{$52.14\%$} & \answer{$20\%$} & \answer{$15.73\%$} & \answer{$100\%$} \\
    \end{tabular}
    \end{center}
\end{remarquebox}

\begin{remarquebox}
    Une fréquence est toujours comprise entre \answer{$0$} et \answer{$1$} \\

    La somme des fréquences est égale à \answer{$1$} ou (ou à \answer{$100\%$})

\end{remarquebox}

